\subsection{4.4.4. 进展性}\label{progress}

把提供方的撤出推迟到其依赖者消失之后的守卫，只有在其最终释放时才能给出定理~63。
登记处各纤维上的一个关系承载了这一论证。

定义 65. 登记处名字上的优先关系定义为

\[
n \prec m := p_{n} \cap d_{m} \neq \emptyset \tag{60}
\]

使得 \(n\) 可以提供 \(m\) 所声明的键。它只读取 \(d\) 和 \(p\)，而由引理~54(5)，这两个字段
随纤维条目的出现而存在，之后不再被写入。

定理~66 和定理~73 都是在 \(\prec\) 无环的假设下建立的；这是假设而非定义本身给出的性
质——对于声明了自身所提供的键的组件，\(n \prec n\) 成立。\(\prec\) 所排序的是两个纤维的
激活而非它们的存续期：\(n \prec m\) 表示 \(n\) 必须首先变为 \(\mathsf{Active}\)，\(m\) 才能如
此；而提供方比其消费方存活更久则是定理~63(2)，一个关于带守卫演算的定理。

纤维的目标视图既响应创建它的纤维，也响应它的提供方。创建者通过定义~47 的原语写入的
是 \(\tau_{n}\)，而由引理~54(5)，\(\tau\) 是单调的。因此创建者在其子纤维的整个存续期内至
多能转动其子纤维的目标视图一次。

进展性断言某条规则可应用，因此它针对宿主必须提供的规则来表述：L-Begin、L-Leave、
L-Unload、着陆规则 L-Iter、L-Finish 和 L-Raise，以及 L-Divert。它从不诉诸 L-Divert 的中止
分支，因此受第 4.3.3 节惯性约束的宿主也被涵盖。

定理 66. （进展性。）假设 \(\prec\) 无环、对每个 \(n\) 有 \(\mathrm{len}(e_{n}) \leq K\)、且定
义~60 中的名字集合 \(N\) 有限；并设每一步都应用某条生命周期规则。记 \(S(n)\) 为作用在
\(n\) 上的步数，记

\[
V(n) := \left| \left\{ t : \mathrm{target}_{n}^{t} \neq \mathrm{target}_{n}^{t+1} \right\} \right| \tag{61}
\]

为其目标视图转动的次数。则

\begin{enumerate}
\def\labelenumi{\arabic{enumi}.}
\item
  （无死锁。）\(\neg \mathrm{quiet}^{t}\) 蕴含在 \(\gamma^{t}\) 处有某条生命周期规则可应用；
\item
  （终止性。）\(S(n) \leq (K+4)(V(n)+1)\)，且 \(V(n)\) 与 \(\sum_{n} S(n)\) 均为有限。
  从而每条极大生命周期步序列都结束于一个静息状态。
\end{enumerate}

证明。无死锁。设 \(\neg \mathrm{quiet}^{t}\) 成立，则某纤维 \(n\) 不满足定义~49 中静息的任一
条款。对照表~1 的四种情形，它可以是：

\begin{itemize}
\item
  \(\theta_{n}^{t} = \mathsf{Inactive}(\bot)\) 且 \(\mathrm{target}_{n}^{t} \neq \bot\)：
  L-Begin 可应用；
\item
  \(\theta_{n}^{t} = \mathsf{Reloading}(-,-,\omega_{n})\) 且 \(\mathrm{target}_{n}^{t} =
  \omega_{n}\)：L-Iter、L-Finish 和 L-Raise 中由 \(i_{n}^{t}(\gamma^{t})\) 的值所选中的那条可
  应用；
\item
  \(\theta_{n}^{t} = \mathsf{Reloading}(-,-,\omega_{n})\) 且 \(\mathrm{target}_{n}^{t} \neq
  \omega_{n}\)：若 \(i_{n}^{t}(\gamma^{t})\) 抛出，则 L-Raise 可应用；否则由 L-Divert 应用，
  令该次迭代着陆而非中止；
\item
  \(\theta_{n}^{t} = \mathsf{Active}(-,\omega_{n})\) 且 \(\mathrm{target}_{n}^{t} \neq \omega_{n}\)：
  L-Leave 可应用。
\end{itemize}

设没有任何纤维属于以上四类，于是剩下某个满足 \(\theta_{m_{0}}^{t} =
\mathsf{Unloading}(-,-,-)\) 的 \(m_{0}\)。按如下方式构造 \(m_{0}, m_{1}, \ldots\)：给定处
于 \(\mathsf{Unloading}\) 的 \(m_{j}\)，要么 \(\neg \mathrm{relied}_{m_{j}}^{t}\)，此时 L-Unload
可作用于 \(m_{j}\)，构造停止；要么存在 \(m_{j+1} \neq m_{j}\) 和 \(k_{j}\)，使得
\(\mathrm{installed}_{m_{j+1}}^{t}\) 且 \(\omega_{m_{j+1}}^{t}(k_{j}) = m_{j}\)。在后一种情形
下

\[
k_{j} \in d_{m_{j+1}} \cap \operatorname{dom}\Big(\sigma_{m_{j}}^{t}\Big) \subseteq d_{m_{j+1}} \cap p_{m_{j}}
\]

第二个成员关系是定理~63(3) 应用于 \(m_{j+1}\) 中 \(t\) 所处的那个片段，从而
\(m_{j} \prec m_{j+1}\)。此外 \(\mathrm{target}_{m_{j+1}}^{t} \neq \omega_{m_{j+1}}^{t}\)：处
于 \(\mathsf{Unloading}\) 的纤维位于定义 \(\sigma_{\gamma}\) 的并集之外，因此在 \(\gamma^{t}\)
处 \(k_{j}\) 要么未被提供，要么由 \(m_{j}\) 之外的纤维提供。若 \(m_{j+1}\) 处于
\(\mathsf{Active}\) 或 \(\mathsf{Reloading}\)，则它将属于已被排除的四类之一，所以它处于
\(\mathsf{Unloading}\)，构造继续。\(m_{j}\) 序列按 \(\prec\) 递增，从而由无环性各不相等，而
\(\operatorname{dom}(F^{t})\) 有限，所以构造停止。

终止性。两条断言界定 \(S(n)\)。

\begin{enumerate}
\def\labelenumi{(\Alph{enumi})}
\item
  在某个使 \(\mathrm{target}_{n}^{t}\) 恒等于 \(\omega^{*}\) 的极大区间上，至多有 \(K+4\) 步
  作用于 \(n\)。读取表~1 的 \(\theta_{n}\) 列：从 \(\omega \neq \omega^{*}\) 的
  \(\mathsf{Active}(-,\omega)\) 出发，纤维执行一次 L-Leave 和一次 L-Unload，然后，若
  \(\omega^{*} \neq \bot\)，执行一次 L-Begin 和至多 \(\mathrm{len}(e_{n}) \leq K\) 次着陆，
  若最后一次着陆是 L-Raise，则还有第二次 L-Unload；从针对 \(\omega \neq \omega^{*}\) 的
  \(\mathsf{Reloading}\) 出发，它用一次 L-Divert 代替 L-Leave；从任何其他状态出发，它执行
  上述序列的一个后缀。在该区间内不会再有 L-Divert 或 L-Leave，因为 L-Begin 所写入的
  \(\omega\) 正是 \(\mathrm{target}_{n}^{t} = \omega^{*}\) 本身；而在 \(\mathsf{Active}(-,\omega^{*})\)
  处、\(\omega^{*} = \bot\) 时的 \(\mathsf{Inactive}(\bot)\) 处以及 \(\mathsf{Inactive}(\xi)\) 处，
  没有任何规则可应用。
\item
  若 \(\mathrm{target}_{n}^{t} \neq \mathrm{target}_{n}^{t+1}\) 且步 \(t\) 作用于 \(m\)，则要
  么 \(m \prec n\)，要么步 \(t\) 写入 \(\tau_{n}\)。由定义~46，target 的值是 \(\tau_{n}\) 与
  \(d_{n}\) 中键的提供方的表的一个函数；提供方满足 \(k \in \operatorname{dom}(\sigma_{m})
  \cap d_{n}\)，从而 \(m \prec n\)，而由引理~54(1)，表只会在作用于其自身纤维的步上变化。
  在第一种情形中无环性给出 \(m \neq n\)，而引理~54(5) 的单调性允许第二种情形在每个纤维
  上至多发生在一个 \(t\) 处。
\end{enumerate}

由 (A)，区间计数将 \(S(n)\) 界定为 \(S(n) \leq (K+4)(V(n)+1)\)；由 (B)，target 的每一次
转动要么消耗严格位于 \(n\) 之下 \(\prec\) 序中的某个纤维的一步，要么就是 \(\tau_{n}\) 所提供
的唯一一次转动，因此 \(V(n) \leq 1 + \sum_{m \prec n} S(m)\)。由于 \(\prec\) 无环且 \(N\) 有
限，递归

\[
B(n) := (K+4)\left( 2 + \sum_{m \prec n} B(m) \right)
\]

是良基的，并定义出满足 \(S(n) \leq B(n)\) 的 \(B\)；从而 \(V(n)\) 有限，且
\(\sum_{n} S(n) \leq \sum_{n} B(n)\)。由 (1)，任何无法继续扩展的序列都是静息的。\(\square\)

\(N\) 的有限性是假设而非推导出来的，组件上的一个条件即可给出它。宿主所持有的组件是在
一切运行之前给定的有限多个程序，因此，若没有任何组件能够——哪怕是间接地——注册一个
注册了自身实例的组件的纤维，那么这些注册构成一棵有界深度的树，而
\(\mathrm{len}(e_{n}) \leq K\) 界定了它的分支。该假设排除的是无界地注册自身实例的组件。

target 记录的是提供纤维而不是一个布尔值；在第 4.2 节的单源约定下，二者驱动相同的转移，
因为那里每个键至多有一个可能的提供方。该视图买来的是上述结果的词汇：定理~63 和定
理~64 都谈到纤维激活时针对的解析，也正是这一点使这些结果在 3.2.3 节的作用域解析下仍
然成立——在那里一个键在不同领域可解析到不同提供方，供给不再强制该视图。实现承载了这
一作用域，并将视图保存在 fiber.committed 中（见 5.1.3 节）。

\subsection{4.4.5. 合流性}\label{confluence}

迄今为止的结果都是关于单个纤维的。刻画系统整体性质的特性是：其动态历史不留下任何痕
迹——无论一个运行中的系统经历过怎样的激活与去激活序列，它静息时所在的状态，正是同
样的插入与退役操作在下列情形下会产生的那个状态：每个最终处于激活状态的组件都按依
赖顺序被加载一次，并且从不被卸载。生命周期关系是合流的，它所收敛到的范式正是静态组
装出来的那个。就动态组合而言，这相当于变更传播为增量计算所确立的、与从头求值相一致
的性质 {[}45{]}。

该断言只关于 \(\longrightarrow\)。编排步是输入，给定不同输入的两条序列落在不同位置并无有
趣的原因；问题在于，生命周期规则——它们在下一步作用于哪个纤维、以及处于
\(\mathsf{Reloading}\) 的纤维走哪个出口上都是非确定的——能否被弄得彼此分歧。

首先需要三条引理。第一条不参照任何步序列就确定最终处于 \(\mathsf{Active}\) 的纤维集合，
这正是使它成为输入而非调度的函数的原因。

定义 67. 当纤维在 \(\gamma\) 处未被退役、注册它的纤维受支撑、且它声明的每个键都由某个
受支撑的纤维提供时，称该纤维在 \(\gamma\) 处受支撑。\(\operatorname{dom}(F_{\gamma})\) 上的支
撑关系是上述两条款所读取的两个关系的并，

\[
m \vartriangleleft n := m \prec n \vee \pi_{n} = m \tag{62}
\]

并且在它良基之处（引理~68），我们把支撑集，即在 \(\gamma\) 处受支撑的纤维，记作 \(A\)：

\[
n \in A := \lnot \tau_{n} \land (\pi_{n} = \mathsf{root} \lor \pi_{n} \in A) \land \forall k \in d_{n}.\, \exists m \in A.\, k \in p_{m} \tag{63}
\]

其中 \(\pi_{n} = \mathsf{root}\) 标记编排器插入的纤维，而 \(\pi_{n}\) 否则指其激活注册了
\(n\) 的纤维。这些条款除 \(\tau, \pi, d, p\) 外不读取任何字段。两个半边都把纤维关联到紧
挨其下的纤维——是父级而非祖先，是直接提供方而非传递提供方，因为条款读取的正是这些；
当后续结果需要某种序时，它们取传递闭包，其极小元、极大元和线性化与
\(\vartriangleleft\) 的相同。

这些条款引用了 \(A\) 本身，因此该定义是沿 \(\vartriangleleft\) 的递归，而正是下面这条使它成
为一个有解的递归。

引理 68. （支撑是良基的。）设 \(\prec\) 无环，且 \(\gamma\) 由一条步序列到达。则
\(\vartriangleleft\) 是良基的，且 \(A\) 是定义~67 的唯一解，一个仅由 \(\tau, \pi, d, p\) 决定
的函数。

证明。按注册每个名字的步的索引来排序 \(\operatorname{dom}(F_{\gamma})\) 中的名字，定义~53
通过让序列从空登记处开始提供了这一索引。\(\vartriangleleft\) 的父级半边沿该索引下降：
O-Insert 以 \(\pi \in \operatorname{dom}\big(F_{\gamma}\big)\) 为前提，因此父指针所指向的纤维注
册得更早，迭代该指针可在有限步内到达一个名字的全部祖先。因此一个环必须使用
\(\prec\)，而由于 \(\prec\) 无环，它必须混合两者，这就需要某个 \(m\) 声明一个由 \(m\) 自身
子树中的某纤维提供的键。这样的纤维由 \(m\) 或 \(m\) 的某个后代的激活注册，因此在
\(m\) 的 L-Begin 之后的一步注册；该 L-Begin 以 \(\gamma \models d_{m}\) 为前提，所以在它之
前提供该键的纤维已是 \(\mathsf{Active}\)，而定义~58 的第 (2) 款使该键不可能有第二个提供
方。因此那个会闭合环的纤维永远不会被注册，这条边不在 \(\operatorname{dom}(F_{\gamma})\) 中。
良基递归有唯一解，且这些条款只读取四个字段。\(\square\)

最后一条读取 \(p\)，即组件可能提供的键；而 target 读取 \(\operatorname{dom}(\sigma_{\gamma})\)，
即其纤维已安装的键，定义~43 仅通过 \(\operatorname{dom}(\sigma_{n}) \subseteq p_{n}\) 将二者关
联。因此支撑集在一般情况下过度近似 \(\mathsf{Active}\) 纤维，而弥合这一差距的条件如下。

定义 69. 当组件的某个完成的激活已安装 \(p\) 中的每一个键，使得在实例化它的每个
\(\mathsf{Active}\) 纤维处 \(\operatorname{dom}(\sigma_{n}) = p_{n}\) 时，称组件 \((d, p, e)\)
在其供给上是完全的。

与独立性（定义~60）一样，这是只作用于组件的条件，不涉及任何生命周期状态或步；而独
立性已经界定了它可能失效到何种程度：如果某个组件只在另一个组件的效应所及的上下文
状态中才安装某个键，那么它的前向映射就无法与那个组件的前向映射交换，因此纤维安装
的键由它的组件而非调度决定。完全性所增加的是：这个确定的集合是全部 \(p\)，而不是它
的真子集。

引理 70. （静息时的支撑。）设 \(\prec\) 无环、\(\mathrm{quiet}(\gamma)\) 成立、\(\gamma\) 中没有
纤维失败，且 \(\gamma\) 的每个组件在其供给上完全（定义~69）。则支撑集就是
\(\mathsf{Active}\) 纤维的集合：

\[
A = \{ n : \theta_{n} = \mathsf{Active}(-,-) \} \tag{64}
\]

证明。记右端为 \(A^{\prime}\)。由于没有纤维失败，定义~49 的静息只留下
\(\mathsf{Inactive}(\bot)\) 和 \(\mathsf{Active}\) 两种状态，并给出

\[
n \in A^{\prime} \Longleftrightarrow \mathrm{target}_{n}(\gamma) \neq \bot
\]

由定义~46，右端恰在 \(\neg \tau_{n}\) 且每个 \(k \in d_{n}\) 都位于
\(\operatorname{dom}(\sigma_{\gamma})\) 中时成立，而由定义~69，
\(\operatorname{dom}(\sigma_{\gamma}) = \bigcup_{m \in A^{\prime}} p_{m}\)。中间那条正是 target
不再承载的条款，由注册来提供：\(\pi_{n} \neq \mathsf{root}\) 的纤维只由 \(\pi_{n}\) 的激活
注册，若 \(\pi_{n} \notin A^{\prime}\)，则 \(\pi_{n}\) 不是 \(\mathsf{Active}\)，于是它的累加器
已经运行，并按定义~47 使 \(n\) 退役，给出 \(\tau_{n}\)。因此 \(A^{\prime}\) 满足定义~67 的
各条款，而引理~68 给出唯一解，所以 \(A = A^{\prime}\)。\(\square\)

引理 71. （换位。）设各步两两独立且 \(F^{t}\) 良构，并设步 \(t\) 和 \(t+1\) 作用于不同纤
维 \(m\) 和 \(n\)。

\begin{enumerate}
\def\labelenumi{\arabic{enumi}.}
\tightlist
\item
  若两者都应用激活规则，即 L-Begin、L-Iter 或 L-Finish，且步 \(t+1\) 在 \(\gamma^{t}\) 处可
  应用，则步 \(t\) 在步 \(t+1\) 从 \(\gamma^{t}\) 产生的状态处可应用，且两种次序到达相同的
  \(\gamma^{t+2}\)；
\item
  若步 \(t\) 在 \(m\) 处应用激活规则、步 \(t+1\) 在 \(n\) 处应用编排规则，且步 \(t\) 不注
  册 \(n\)，则二者同样成立。
\end{enumerate}

证明。对 (1)，由表~1，\(m\) 的步写入 \(\theta_{m}\)，并在 \(\Psi^{t} \in \mathfrak{M}(e_{m})\)
之内写入表 \(\sigma_{m}\) 和效应部分。因此它不碰 \(\theta_{n}\) 和 \(i_{n}\)，且由定义~60 的
第二条条件，也不碰 \(i_{n}\) 所产出的逆与续体，所以只需检查步 \(t+1\) 中提及
\(\mathrm{target}_{n}\) 的前提。其退役半边不会失效，因为没有激活规则写入 \(\tau\)。其解析
半边也不能变动：步 \(t+1\) 在 \(\gamma^{t}\) 处可应用使每个 \(k \in d_{n}\) 都位于
\(\operatorname{dom}(\sigma^{t})\) 中，而定义~58 的第 (2) 款使提供这样的 \(k\) 的纤维成为唯
一可能的提供者，因此 \(k \notin p_{m}\)，对 \(\sigma_{m}\) 的任何写入都不会触及 \(d_{n}\) 中
的键。反向的同样论证使步 \(t\) 保持可应用。最后，由定义~60 的第一条条件，
\(\Psi^{t} \in \mathfrak{M}(e_{m})\) 与 \(\Psi^{t+1} \in \mathfrak{M}(e_{n})\) 交换，且两次编辑
写入的是不同纤维的控制字段，所以无论哪种次序，复合结果都相同。

对 (2)，由表~1，编排步有 \(\Psi^{t+1} = \mathrm{id}_{\Gamma}\)，因此两个状态映射完全交换；
它的编辑 \(\mathrm{edit}^{t+1}\) 只在 \(n\) 处写入 \(\tau_{n}\) 或
\(\operatorname{dom}(F_{\gamma})\)，激活步既不读取也不写入这些：激活步的前提读取
\(\theta_{m}, i_{m}, \tau_{m}\) 和 \(\mathrm{target}_{m}\)，而插入新 \(n\) 的 O-Insert 不移动任
何 target——新纤维不提供任何东西；而对 \(n\) 的 O-Retire 或 O-Remove 使 \(\sigma_{\gamma}\)
保持原样——在一种情形下 \(n\) 处于 \(\mathsf{Inactive}\)，在另一种情形下它的表不受影响。所
以步 \(t\) 保持可应用。反之，编排步的每个前提要么在 \(n\) 处读取（步 \(t\) 不写入那里），
要么是 O-Insert 的两个前提之一——更小的登记处只会放松这些前提——因此它在
\(\gamma^{t+1}\) 处的可应用性蕴含它在 \(\gamma^{t}\) 处的可应用性；这里正是步 \(t\) 不注册
\(n\)，才使 \(n\) 在 O-Retire 和 O-Remove 要求它出现之处仍存在于 \(\gamma^{t}\)。\(\square\)

引理 72. （删除。）设步序列两两独立、每个组件在其供给上完全（定义~69）、它到达一个
无纤维失败的静息状态 \(\gamma^{T}\)、\([b, u]\) 是 \(n\) 的一个闭合的片段、序列中没有满足
\(n \prec m\) 的 \(m\) 的片段闭合、且 \(n\) 在 \([b, u]\) 期间注册的任何纤维都没有片段。记
这些注册所抽取的名字为 \(R\)。那么，删除 \([b, u]\) 中作用于 \(n\) 的步，连同作用于 \(R\)
中某个名字的每一步，得到一条步序列，它到达的状态与 \(\gamma^{T}\) 在 \(\approx\) 意义下相等，
且在 \(R\) 之外与它 \(\simeq\) 意义下相等。

证明。被删除的步让状态保持原样。设 \(t_{1} < \cdots < t_{l}\) 是 \([b, u]\) 中作用于 \(n\) 之
外纤维的步。推论~62 给出

\[
\gamma^{u+1} \approx \bigl( \Psi^{t_{l}} \circ \cdots \circ \Psi^{t_{1}} \bigr) \bigl( \gamma^{b} \bigr)
\]

其右端正是 \([b, u]\) 中幸存的步自行产生的结果：\(\gamma^{b-1} \approx \gamma^{b}\)，且它们的
编辑写入的是删除不触及的、\(n\) 之外纤维的控制字段。由表~1，\(n\) 的被删除步不写入
\(\theta_{n}\) 之外的任何字段，而引理~54(4) 在 \(u\) 处把它恢复到 \(\mathsf{Inactive}(\bot)\)
（没有纤维失败），且它在 \(\gamma^{b-1}\) 处正是这个值。

一个不变量承载着后缀。把幸存步在对应于点 \(t\) 处到达的状态记为 \(\gamma^{\prime t}\)。我
们断言：对每个 \(t > u\)，\(\gamma^{t} \approx \gamma^{\prime t}\)；\(R\) 中的每个名字在
\(\gamma^{t}\) 处都是残留的、在 \(\gamma^{\prime t}\) 中不存在；两个状态在 \(R\) 之外的每个
名字的每个字段上都一致。在 \(t = u+1\) 处，这由上一段连同定义~47 给出：定义~47 使
\(R\) 中的每个名字被在 \(u\) 处运行的累加器退役、处于 \(\mathsf{Inactive}(\bot)\) 且持有空
表——由假设，\(R\) 的纤维没有片段。归纳步是对 \(R\) 中每个名字依次应用引理~57(1)：作
用于 \(R\) 之外的步在两个状态处有相同的前提，到达再次满足 \(\approx\) 的状态，并使 \(R\)
的条目保持残留。作用于 \(R\) 中某个名字的步属于被删除者，而引理~57(2) 说明了为什么必
须删除而非保留——对不存在名字的 O-Retire 或 O-Remove 没有纤维可作用；同样由 (1)，这
样的步不移动 \(R\) 之外的任何字段，所以丢弃它保持不变量。因此最终状态满足 \(\approx\)，
且在 \(R\) 之外相等。

没有任何幸存步失去前提。作用于 \(m \notin R \cup \{n\}\) 的步只通过
\(\mathrm{target}_{m}(\gamma)\) 或 \(\mathrm{relied}_{m}(\gamma)\) 读取 \(n\)。前者在 \(m\) 声
明 \(n\) 提供的键时依赖 \(n\)，从而 \(n \prec m\)；也在 \(n\) 注册了 \(m\) 时依赖 \(n\)，这
使 \(m \in R\)。在第一种情形下，由假设，\(m\) 的片段不闭合，所以它在 \(\gamma^{T}\) 处是打
开的；在那里静息给出 \(\omega_{m} = \mathrm{target}_{m}^{T}\)，而引理~70 把它的值归入
\(\mathsf{Active}\) 纤维——而 \(n\) 不是 \(\mathsf{Active}\)；由于一个键至多有一个可能的提供
方，\(n\) 在 \(m\) 的 L-Begin 时也没有提供 \(d_{m}\) 中的任何键。后者只通过 \(\omega_{n}\) 的
值读取 \(n\)，而删除该片段只会使 \(\mathrm{relied}\) 变为假，这放宽了 L-Unload 上的守卫而
不是阻塞它。这样的步对 \(R\) 中某个名字的读取由不变量涵盖。两两独立性是效应函数的性
质，所以删除步保持它。\(\square\)

定理 73. （合流性。）设一条步序列到达一个无纤维失败的静息状态 \(\gamma^{T}\)，各步两两独
立、每个组件在其供给上完全（定义~69），且 \(A\) 如定义~67 中那样。则

\begin{enumerate}
\def\labelenumi{\arabic{enumi}.}
\item
  （规范形。）\(\gamma^{T}\) 可以从 \(\gamma^{0}\) 出发由一条序列到达——除归约所撤回其条目的
  那些名字不计——该序列按原顺序执行相同的编排步：作用于编排器插入的纤维的编排步位
  于每个生命周期步之前，其余每个编排步位于注册其作用纤维的那一步之后；并且，对
  \(A\) 的一个使 \(\vartriangleleft\) 线性化的枚举 \(n_{1}, \ldots, n_{k}\)，按该顺序取每个
  \(n_{i}\) 的一个片段。
\item
  （合流性。）任何两条从 \(\gamma^{0}\) 出发、执行相同编排步的此类序列，经过如引理~56 那
  样的重命名后，到达由 \(\simeq\) 和 \(\approx\) 关联的状态。
\end{enumerate}

证明。对 (1)，序列的片段有两类：闭合的，以及在 \(\gamma^{T}\) 处仍打开的——后者由静息
\(\mathrm{quiet}^{T}\) 和引理~70 可知是 \(A\) 中每个纤维各一个片段。

闭合片段优先处理，对其个数做归纳。每个阶段在片段仍闭合的纤维中选取一个在
\(\vartriangleleft\) 序下极大的纤维 \(n\) 的闭合片段；由引理~68 和 \(N\) 的有限性，这样的片段
存在。于是引理~72 的三个假设都满足。由极大性，没有任何满足 \(n \prec m\) 的 \(m\) 有闭
合片段。而 \(n\) 在 \([b, u]\) 期间注册的纤维也没有片段：这样的纤维被在 \(u\) 处运行的累
加器退役（定义~47），并由引理~54(5) 保持退役，所以它的目标视图为 \(\bot\)，引理~70 把
它置于 \(A\) 之外，从而它在 \(\gamma^{T}\) 处没有打开的片段；且 \(\vartriangleleft\) 通过它的父
指针把它与 \(n\) 关联，所以由极大性它也没有闭合片段。引理移除了该片段以及它所注册名
字的各步，使 \(\gamma^{T}\) 除那些名字之外保持原样。测度减一，所以不再有闭合片段。

\(A\) 之外的纤维不取任何生命周期步。由引理~70 和静息 \(\mathrm{quiet}^{T}\)，它在
\(\gamma^{T}\) 处没有打开的片段，而现在也不再有闭合片段，所以它根本没有片段，并始终处
于 \(\mathsf{Inactive}(\bot)\)；L-Begin 是那里唯一可应用的规则，而应用它会打开一个片段。

接下来处理编排步。作用于编排器插入的纤维的编排步，由引理~71(2) 可越过另一纤维的生
命周期步向前移动一位；该引理适用是因为 \(A\) 中某纤维的步不会注册这样的名字：注册抽
取的都是新名字，而这里的名字是原始序列中某次 O-Insert 引入的。与同一纤维的生命周期
步之间则无可交换：\(n\) 的 O-Insert 已经位于 \(n\) 的每一步之前，而 \(n\) 的 O-Retire 或
O-Remove 只作用于 \(A\) 之外，后者不取生命周期步。依次把每个编排步移到最前保持它们
的相对顺序。作用于某次激活所注册的纤维的编排步不能移到最前——它的前提要求该纤维
在场——所以它停留在注册所放的位置；由上一段它作用于 \(A\) 之外，因此由引理~71 的同
一款，它与自身和注册之间的所有步都交换。

片段被排序并连续化，对 \(|A|\) 做归纳。设 \(n_{1}\) 在 \(A\) 中 \(\vartriangleleft\)-极小。则
\(d_{n_{1}} = \emptyset\) 且 \(\pi_{n_{1}} = \mathsf{root}\)，因为定义~67 把 \(d_{n_{1}}\) 中键
的一个提供方和注册 \(n_{1}\) 的纤维放入 \(A\)，而 \(\vartriangleleft\) 把二者都放在 \(n_{1}\) 之
下。所以 \(\mathrm{target}_{n_{1}}\) 不读取另一纤维的任何字段，并且——不再有编排步来写
\(\tau_{n_{1}}\)，也没有 \(n_{1}\) 之下的纤维来使它退役——它恒为常数。作用于 \(n_{1}\) 的每
一步都是激活步，没有片段闭合，且其剩余前提读取 \(\theta_{n_{1}}\) 和 \(i_{n_{1}}\)——由表~1，
只有 \(n_{1}\) 写入这些；因此每一步在每个更早的状态处都可应用，而引理~71 在不移动端点
的前提下把它向前移一位。每次应用都使 \(n_{1}\) 的一步之前的其他纤维步数减一，所以
\(n_{1}\) 的片段变成一个初始的连续块。该论证在紧随该块之后的后缀上对
\(A \setminus \{n_{1}\}\) 重复：在那里 \(n_{1}\) 全程处于 \(\mathsf{Active}\) 且不再取步，所以它
也贡献一个常数 target。这样产生的枚举按构造使 \(\vartriangleleft\) 线性化。

对 (2)，两条序列都通过 (1) 归约到某条规范序列，而两次归约在同一个（差一个重命名的）
\(A\) 上进行。定义~67 读取 \(\tau, \pi, d, p\)，其中后三个字段随纤维条目的出现而只写一次
（引理~54(5)），因此需要看清的是：相同的名字带着相同的 \(d, p, \pi\) 出现，且相同的名
字被退役。插入操作按假设为两条序列所共有。注册操作也为二者共有：\(A\) 中某纤维的一
次激活在其每次迭代中注册迭代器在那里点名的组件，而定义~60 的第二条条件使这一点在
交错中保持不变，所以 \(A\)-纤维之下的注册树是该纤维组件的函数；这些注册所抽取的名
字并不共有，正是在这里应用引理~56，用双射匹配两棵树。而一次退役要么是编排步（共
有），要么是累加器执行的 O-Retire，后者退役的恰好是同一次激活所注册的名字。使
\(\vartriangleleft\) 线性化的两个枚举只相差对不可比片段的换位，而引理~71 再次使端点保持不
变，所以两条规范序列一致。结合定理~66 的终止性，生命周期关系因此具有唯一的范式。
\(\square\)

失败被排除在陈述之外，因为它是分歧的真实来源，而且不应把这套演算解读为否认失败：
一步是否抛出取决于它运行时所处的状态，因此一个调度可能使某纤维失败，而另一个调度
完成它，两个静息状态于是会在该纤维的生命周期状态上不同。由推论~62，它们在除此之外
的任何方面都没有差异——该推论把失败纤维对状态的贡献置为无。

在第 4.2 节的基础演算中同样的定理成立，证明除了去掉一条款外无需任何替换。那里的
L-Unload 不带守卫，所以引理~72 的最后一段是空洞的；该引理的其余部分只诉诸静息
\(\mathrm{quiet}^{T}\)，基础演算原样提供之。

正是这条定理使我们可以像对待静态组装的应用那样对 Cordis 应用进行推理。一个添加组件、
移除组件、替换提供方并撤销替换的编排器，保证会到达它在一开始就把最终组合写下来所
得的状态；而组件作者在推理哪些余效应在作用域内时，可以只针对静息状态进行推理。它
也界定了这一保证的范围：它说的是状态，而不是系统沿途产生的排放，这正是第 6.1 节在
边界内部被追踪的获取与跨越边界的排放之间所作的区分。
